\section*{Chapter \ref{ch:g10:storms-and-cyclones} --- Storms and Tropical Cyclones}

\begin{solution}{exo:g10:storms-and-cyclones:1}
A \omterm{def:g10:storms-and-cyclones:thunderstorm}{thunderstorm} is a storm produced by a cumulonimbus, with \omterm{def:g10:storms-and-cyclones:lightning}{lightning} and
thunder, usually heavy showers and gusts, sometimes hail. \omterm{def:g10:storms-and-cyclones:lightning}{Lightning} is a giant
electric spark that discharges charged regions of the storm. Thunder is the
sound of the air of the \omterm{def:g10:storms-and-cyclones:lightning}{lightning} channel, heated so suddenly that it expands
with a shock.
\end{solution}

\begin{solution}{exo:g10:storms-and-cyclones:2}
A \omterm{def:g10:storms-and-cyclones:tropical-cyclone}{tropical cyclone} is a storm hundreds of kilometres across, formed over a
warm tropical sea around a centre of low pressure, with winds spiralling in
and bands of \omterm{def:g10:storms-and-cyclones:thunderstorm}{thunderstorms}. The \omterm{def:g10:storms-and-cyclones:tropical-cyclone}{eye} is its calm, nearly cloudless centre; the
\omterm{def:g10:storms-and-cyclones:tropical-cyclone}{eyewall} is the ring of towering \omterm{def:g10:storms-and-cyclones:thunderstorm}{thunderstorms} around the \omterm{def:g10:storms-and-cyclones:tropical-cyclone}{eye}. The winds are
strongest in the \omterm{def:g10:storms-and-cyclones:tropical-cyclone}{eyewall}.
\end{solution}

\begin{solution}{exo:g10:storms-and-cyclones:3}
\qty{70}{knots}: category 1. \qty{120}{knots}: category 4.
\qty{140}{knots}: category 5. The last two are major hurricanes (category 3
or more).
\end{solution}

\begin{solution}{exo:g10:storms-and-cyclones:4}
$343 \times 9 \approx \qty{3.1}{km}$, then $343 \times 15 \approx
\qty{5.1}{km}$. The second flash is farther away: the storm is moving away.
\end{solution}

\begin{solution}{exo:g10:storms-and-cyclones:5}
\omterm{def:g10:storms-and-cyclones:lightning}{Lightning} tends to strike the highest point around: a tree standing alone in
a field is that point, and the current of a strike runs down it and spreads
through the ground around it. Go into a building or a closed car, and stay
there until well after the last thunder.
\end{solution}

\begin{solution}{exo:g10:storms-and-cyclones:6}
One knot is $1852/3600 \approx \qty{0.514}{m/s}$. \qty{64}{knots}: $64 \times
0.514 \approx \qty{32.9}{m/s}$, and $32.9 \times 3.6 \approx
\qty{119}{km/h}$. \qty{137}{knots}: $\qty{70.5}{m/s}$, about
\qty{254}{km/h} (the official table, rounded through another unit, gives
\qty{252}{km/h}).
\end{solution}

\begin{solution}{exo:g10:storms-and-cyclones:7}
$44 \times 86\,400 \approx \num{3.8e6}$ flashes a day, and $3.8\times10^6
\times 365.25 \approx \num{1.4e9}$, about 1.4 billion a year.
\end{solution}

\begin{solution}{exo:g10:storms-and-cyclones:8}
$\Delta p = 1010 - 935 = \qty{75}{hPa}$: a bulge of about \qty{75}{cm}.
\end{solution}

\begin{solution}{exo:g10:storms-and-cyclones:9}
They form over the warm tropical oceans, mostly between \qty{10}{\degree} and
\qty{30}{\degree} of latitude in both hemispheres, with very few in the South
Atlantic. They drift west, then curve away towards the poles, where they
weaken over colder water. They are absent from the equator because there the
Earth's rotation gives the inflowing air no turn: without it, the air flows
straight into the low pressure and fills it, and no rotating storm can start.
\end{solution}

\begin{solution}{exo:g10:storms-and-cyclones:10}
$1031 - 1.56 \times 70 \approx \qty{922}{hPa}$; $1031 - 1.56 \times 40
\approx \qty{969}{hPa}$.
\end{solution}

\begin{solution}{exo:g10:storms-and-cyclones:11}
Over land the storm loses its supply of evaporating water, and with it the
\omterm{def:g10:earth-energy-budget:latent-heat}{latent heat} that drives it; the rough surface also slows the winds. Over cold
water evaporation is much weaker, since cold air holds far less water vapour
at saturation: the supply of \omterm{def:g10:earth-energy-budget:latent-heat}{latent heat} dwindles.
\end{solution}

\begin{solution}{exo:g10:storms-and-cyclones:12}
At \qty{0}{\degreeCelsius}: $\sqrt{1.4 \times 8.314 \times 273.15/0.02897}
\approx \qty{331}{m/s}$; at \qty{30}{\degreeCelsius}, with $T =
\qty{303.15}{K}$: \qty{349}{m/s}. For \qty{9}{s}: $9 \times 331 \approx
\qty{2980}{m}$ and $9 \times 349 \approx \qty{3140}{m}$, a difference of
about \qty{160}{m}, about \qty{5}{\percent}. The rule gives \qty{3}{km}:
good enough to know how close the storm is, which is all it is for.
\end{solution}

\begin{solution}{exo:g10:storms-and-cyclones:13}
$\Delta p = 1010 - 940 = \qty{70}{hPa}$: the bulge is about \qty{0.7}{m}, the
wind's part $4.0 - 0.7 = \qty{3.3}{m}$, more than four times as much. At high
tide the water reaches $4.0 + 1.5 = \qty{5.5}{m}$ above mean sea level.
\end{solution}

\begin{solution}{exo:g10:storms-and-cyclones:14}
\omterm{def:g10:storms-and-cyclones:tornado}{Tornado}: one turn is $2\pi \times 100 \approx \qty{628}{m}$, travelled in
$628/90 \approx \qty{7}{s}$. Hurricane: $2\pi \times 30\,000 \approx
\qty{188}{km}$, travelled in $188\,000/60 \approx \qty{3140}{s}$, about
\qty{52}{min}. The \omterm{def:g10:storms-and-cyclones:tornado}{tornado} turns some 450 times faster: it is a small, very
fast whirl, the hurricane a huge and slower one.
\end{solution}

\begin{solution}{exo:g10:storms-and-cyclones:15}
$24/3888 \approx \qty{0.6}{\percent}$; $414/3888 \approx
\qty{10.6}{\percent}$; $121/3888 \approx \qty{3.1}{\percent}$. If latitude
did not matter, about \qty{9}{\percent} of the storms, some 350, would start
within \qty{5}{\degree} of the equator; only 24 do, about fifteen times
fewer. Near the equator the Earth's rotation is too weak to start the spin.
\end{solution}

\begin{solution}{pb:g10:storms-and-cyclones:1}
\textbf{1.} $\pi \times (665\,000)^2 \approx \qty{1.39e12}{m^2}$.

\textbf{2.} $0.015 \times 1.39\times10^{12} \approx \qty{2.08e10}{m^3}$.

\textbf{3.} One cubic kilometre is \qty{e9}{m^3}: about \qty{21}{km^3}.

\textbf{4.} $2.08\times10^{10} \times 1000 \approx \qty{2.08e13}{kg}$.

\textbf{5.} $2.08\times10^{13} \times 2.5\times10^6 \approx \qty{5.2e19}{J}$.

\textbf{6.} $5.2\times10^{19}/86\,400 \approx \qty{6.0e14}{W}$.

\textbf{7.} $6.0\times10^{14}/1.22\times10^{17} \approx 0.005$, about
\qty{0.5}{\percent}.

\textbf{8.} $240 \times 1.39\times10^{12} \approx \qty{3.3e14}{W}$. The storm
releases nearly twice as much heat as the sunlight absorbed by its whole area.

\textbf{9.} The storm does not run on the sunlight that falls on it: it
releases heat stored as water vapour, evaporated from a wider area of warm
sea and drawn in by the winds. The storm concentrates energy gathered over a
much larger region.

\textbf{10.} The same mass as the rain, about \qty{2.1e13}{kg} a day (the
vapour carried in must be renewed as fast as it rains out).

\textbf{11.} $6.0\times10^{14}/82 \approx \qty{7.3e12}{m^2}$, about five times
the area of the storm: the storm must draw vapour from a much wider area, and
its winds make the sea beneath it evaporate far faster than the average.

\textbf{12.} \qty{24.1}{g/kg} at \qty{28}{\degreeCelsius} against
\qty{14.7}{g/kg} at \qty{20}{\degreeCelsius}: $24.1/14.7 \approx 1.64$, about
\qty{64}{\percent} more.

\textbf{13.} Warm water evaporates fast and warm air holds much more vapour:
the inflowing air carries much more \omterm{def:g10:earth-energy-budget:latent-heat}{latent heat} to release, and, saturated,
it rises freely. Over cooler seas the supply is too small to drive the
engine.

\textbf{14.} Over land there is no sea to evaporate: the supply of \omterm{def:g10:earth-energy-budget:latent-heat}{latent
heat} stops, and the rough ground slows the winds.

\textbf{15.} $1.5\times10^{12} \times 86\,400 \approx \qty{1.3e17}{J}$.

\textbf{16.} $96 \times 1852/3600 \approx \qty{49.4}{m/s}$.

\textbf{17.} $\frac12 \times 1.16 \times 49.4^2 \approx \qty{1.4e3}{J}$.

\textbf{18.} $5.2\times10^{19}/1.3\times10^{17} = 400$: the \omterm{def:g10:earth-energy-budget:latent-heat}{latent heat} is
some 400 times the energy given to the winds.

\textbf{19.} Most of it warms the upper air of the storm, which spreads out at
the top and gives the heat off to space as infrared radiation; the winds
themselves lose their energy to friction against the sea and the ground.

\textbf{20.} $1/400 = \qty{0.25}{\percent}$ (the script gives a ratio of
402): \textbf{about one part in 400 of the heat released by condensation
becomes wind}.
\end{solution}
